\subsection{Theta-stratifications}%
\label{sub:Theta-stratifications}

The framework of \(\Theta\)-stratifications and numerical invariants used here is developed in~\cite{HLinstability}.

Let \(\mf{X}\) be an algebraic stack and consider \(\Theta\coloneqq[\A^1/\G_m]\), where the action on points is \(\lambda\cdot z=\lambda z\); equivalently, the coordinate function \(t\) has weight \(-1\). Assume that the mapping stacks below are algebraic. Then define the stack \(\ms{Filt}(\mf{X})\coloneqq\ms{Map}(\Theta,\mf{X})\) of filtered objects and the stack \(\ms{Grad}(\mf{X})\coloneqq\ms{Map}(B\G_m,\mf{X})\) of graded objects. There are canonical morphisms
\[
    \ev_1\colon\ms{Filt}(\mf{X})\to\mf{X},
    \qquad
    \on{gr}\colon\ms{Filt}(\mf{X})\to\ms{Grad}(\mf{X}),
    \qquad
    \sigma\colon\ms{Grad}(\mf{X})\to\ms{Filt}(\mf{X}).
\]
Here \(\ev_1\) is restriction to \(1\in\A^1\), \(\on{gr}\) is induced by the inclusion \(B\G_m=\{0\}/\G_m\hookrightarrow\Theta\), and \(\sigma\) is induced by the projection \(\Theta\to B\G_m\). In particular, \(\on{gr}\circ\sigma=\on{id}\). We also write \(\on{tot}=\ev_1\circ\sigma\colon\ms{Grad}(\mf{X})\to\mf{X}\).

\begin{defn}
    A \textit{\(\Theta\)-stratum} in an algebraic stack \(\mf{Y}\) is a union of connected components \(S\subset\ms{Filt}(\mf{Y})\) such that the restriction
    \[
        \ev_1|_S\colon S\longrightarrow\mf{Y}
    \]
    is a closed immersion.

    A \textit{\(\Theta\)-stratification} of \(\mf{X}\) consists of
    \begin{enumerate}
        \item A totally ordered set \(I\) with a least element \(0\);
        \item An increasing collection of open substacks \(\mf{X}_{\leq\alpha}\subset\mf{X}\), indexed by \(\alpha\in I\), such that
        \[
            \mf{X}_{\leq\alpha}\subset\mf{X}_{\leq\beta}
            \quad\text{for }\alpha\leq\beta,
            \qquad
            \bigcup_{\alpha\in I}\mf{X}_{\leq\alpha}=\mf{X},
        \]
        and such that for every \(x\in|\mf{X}|\), the set \(\{\alpha\in I\mid x\in|\mf{X}_{\leq\alpha}|\}\) has a least element;
        \item For every \(\alpha>0\), a \(\Theta\)-stratum \(S_\alpha\subset\ms{Filt}(\mf{X}_{\leq\alpha})\) such that
        \[
            \mf{X}_{<\alpha}
            \coloneqq
            \bigcup_{\beta<\alpha}\mf{X}_{\leq\beta}
            =
            \mf{X}_{\leq\alpha}\setminus\ev_1(S_\alpha).
        \]
    \end{enumerate}
    The semistable locus is \(\mf{X}^{\mr{ss}}\coloneqq\mf{X}_{\leq0}\), and the stratification gives a disjoint decomposition
    \[
        \mf{X}
        =
        \mf{X}^{\mr{ss}}
        \sqcup
        \bigsqcup_{\alpha>0}\ev_1(S_\alpha).
    \]
\end{defn}

Consider the stack \(\ms{Bun}_G(C) = \ms{Map}(C, BG)\). We will consider the stack
\[ \ms{Map}(\Theta, \ms{Bun}_G(C)) = \ms{Map}(\Theta \times C, BG). \]
Any such map \(f\) is equivalent to a \(\G_m\)-equivariant principal \(G\)-bundle \(\mc{E}\) over \(\A^1\times C\). When \(G=\on{GL}_n\), this is equivalent to a \(\G_m\)-equivariant vector bundle on \(\A^1\times C\), or equivalently to a finitely generated graded locally free \(\mc{O}_C[t]\)-module
\[
    \mc{E}=\bigoplus_{w\in\Z}\mc{E}_w.
\]
With the convention above, the coordinate function \(t\) has weight \(-1\), so multiplication by \(t\) gives maps
\[
    t\colon\mc{E}_w\longrightarrow\mc{E}_{w-1}.
\]
Since \(\mc{E}\) is locally free over \(\mc{O}_C[t]\), multiplication by \(t\) is injective.

\begin{rmk}
    Consider the vector bundle \(E=\mc{E}|_{\{1\}\times C}\). Since \(\G_m\) acts freely on \(\G_m\times C\), equivariant descent gives
    \[
        \mc{E}|_{\G_m\times C}\simeq P_C^*E,
    \]
    where \(P_C\colon\G_m\times C\to C\) is the projection. Equivalently,
    \[
        E\simeq\big(\mc{E}[t^{-1}]\big)_0.
    \]
    Define
    \[
        F^wE\coloneqq t^w\mc{E}_w\subset\big(\mc{E}[t^{-1}]\big)_0=E.
    \]
    Since \(t\mc{E}_{w+1}\subset\mc{E}_w\), we have \(F^{w+1}E\subset F^wE\). The finite generation and local freeness of \(\mc{E}\) imply that this is a finite, exhaustive, and separated decreasing filtration by subbundles, which may be written as
    \[
        0=F^{b+1}E\subset F^bE\subset\cdots\subset F^aE=E.
    \]

    In the other direction, let \(E\) be a vector bundle with a finite decreasing filtration \(F^\bullet E\) by subbundles such that the quotients \(F^wE/F^{w+1}E\) are locally free. The Rees construction is the graded \(\mc{O}_C[t]\)-module
    \[
        \ms{Rees}(E,F^\bullet)
        \coloneqq
        \bigoplus_{w\in\Z}t^{-w}F^wE
        \subset
        E\otimes_{\C}\C[t,t^{-1}].
    \]
    It is locally free over \(\mc{O}_C[t]\), and its fibers at \(t=1\) and \(t=0\) satisfy
    \[
        \ms{Rees}(E,F^\bullet)|_{t=1}\simeq E,
        \qquad
        \ms{Rees}(E,F^\bullet)|_{t=0}
        \simeq
        \bigoplus_w F^wE/F^{w+1}E.
    \]
\end{rmk}

Therefore, a map from \(\Theta\) to the stack of vector bundles is equivalent to a finite \(\Z\)-weighted filtration of a vector bundle by subbundles with locally free associated graded pieces.


\subsection{Non-abelian localization}%
\label{sub:Non-abelian localization}

We will begin with the classical setting. Let \(X\) be a smooth projective variety with an action of \(\G_m\), let \(F\) be a \(\G_m\)-equivariant perfect complex, and write
\[
    X^{\G_m}=\bigsqcup_\alpha Z_\alpha
\]
for the decomposition of the fixed locus into connected components. After localizing the representation ring \(R(\G_m)\) so that the equivariant Euler classes below become invertible, we have
\[
    \chi_{\G_m}(X,F)
    =
    \sum_\alpha
    \chi_{\G_m}\ab(
        Z_\alpha,
        F|_{Z_\alpha}\otimes e(N_{Z_\alpha/X})^{-1}
    ),
\]
where
\[
    e(N_{Z_\alpha/X})
    =
    \lambda_{-1}(N_{Z_\alpha/X}^\vee)
    =
    \sum_{i\geq0}(-1)^i[\Lambda^iN_{Z_\alpha/X}^\vee].
\]

We will now consider the stacky version. Let
\[
    \mf{X}
    =
    \mf{X}^{\mr{ss}}
    \sqcup
    \bigsqcup_\alpha\ev_1(S_\alpha)
\]
be a \(\Theta\)-stratification. Since \(\ev_1|_{S_\alpha}\) is a closed immersion, we may identify \(S_\alpha\) with its image in \(\mf{X}\). The center of the stratum is
\[
    Z_\alpha\coloneqq\sigma^{-1}(S_\alpha)\subset\ms{Grad}(\mf{X}),
\]
where \(\sigma\colon\ms{Grad}(\mf{X})\to\ms{Filt}(\mf{X})\) is induced by the projection \(\Theta\to B\G_m\).

\begin{thm}[Virtual non-abelian localization {\cite[Theorem~5.1]{HLlocalization}}]
    Let \(\mf{X}\) be a quasi-smooth derived stack with a \(\Theta\)-stratification
    \[
        \mf{X}
        =
        \mf{X}^{\mr{ss}}
        \sqcup
        \bigsqcup_\alpha S_\alpha,
    \]
    where we identify every stratum with its image under \(\ev_1\), and let \(Z_\alpha\) be its center. Decompose the restriction of the cotangent complex according to the \(\G_m\)-weights as
    \[
        \L_{\mf{X}}|_{Z_\alpha}
        \simeq
        \L_\alpha^-
        \oplus
        \L_\alpha^0
        \oplus
        \L_\alpha^+,
    \]
    where the homology of \(\L_\alpha^-\), \(\L_\alpha^0\), and \(\L_\alpha^+\) is concentrated in negative, zero, and positive weights, respectively. Define
    \[
        \E_\alpha
        \coloneqq
        \on{Sym}\ab(
            \L_\alpha^-
            \oplus
            (\L_\alpha^+)^\vee
        )
        \otimes
        \det(\L_\alpha^+)^\vee
        [-\on{rk}(\L_\alpha^+)]
        \in
        \ms{QCoh}(Z_\alpha).
    \]
    Here \(\det(\L_\alpha^+)\) and \(\on{rk}(\L_\alpha^+)\) denote the determinant and virtual rank of the perfect complex \(\L_\alpha^+\). We say that a complex of vector spaces is finite-dimensional if the direct sum of all of its cohomology groups is finite-dimensional.
    Let \(F\in\ms{Perf}(\mf{X})\). Assume that
    \[
        R\Gamma\ab(
            Z_\alpha,
            F|_{Z_\alpha}\otimes\E_\alpha
        )
    \]
    is finite-dimensional for every \(\alpha\) and is acyclic for all but finitely many \(\alpha\). Then \(R\Gamma(\mf{X},F)\) is finite-dimensional if and only if \(R\Gamma(\mf{X}^{\mr{ss}},F)\) is finite-dimensional, and
    \[
        \chi(\mf{X},F)
        =
        \chi(\mf{X}^{\mr{ss}},F)
        +
        \sum_\alpha
        \chi\ab(
            Z_\alpha,
            F|_{Z_\alpha}\otimes\E_\alpha
        ).
    \]
    Here quasi-smooth means that \(\L_{\mf{X}}\) is perfect of Tor-amplitude contained in \([-1,1]\).
\end{thm}


\subsection{Geometric invariant theory}%
\label{sub:Geometric invariant theory}

For geometric invariant theory, see~\cite{MFK1994}. The vector-bundle example below uses the Harder--Narasimhan filtration~\cite{HN1975}.

Let \(G\) be a reductive group and \(X\) be a projective scheme equipped with a \(G\)-equivariant line bundle \(\mc{L}\). The semistable locus can be defined using the Hilbert-Mumford criterion, where it consists of all points where for all cocharacters
\(\lambda \colon \G_m \to G\), we can define the point 
\[ x_0 = \lim_{t \to 0} \lambda(t) x \]
and define the quantity
\[ \mu^L(x, \lambda) \coloneqq \frac{-\on{wt}_{x_0}(\mc{L})}{\norm{\lambda}}. \]
A point is semistable if \(\mu^L(x, \lambda) \leq 0\) for all cocharacters of \(G\). The stable locus consists of all points where the inequality is strict and the stabilizer is finite.

\begin{thm}[{\cite{MFK1994}}]
    The stable and semistable loci are open in \(X\). Moreover, \(X^{\on{ss}}/G\) is a good quotient and \(X^s/G\) is a geometric quotient.  
\end{thm}

\begin{exm}
    Consider the moduli stack \(\ms{Bun}_{r,d}\) of vector bundles of rank \(r\) and degree \(d\) on a curve of genus at least \(2\). Then recall that the semistable locus was defined to consist of slope semistable vector bundles. In fact, this can be obtained via GIT.

    Fix an ample line bundle \(\mc{O}(1)\) on \(C\). Then there exists some \(n_0\) such that for all \(n \geq n_0\) and \(E \in \ms{Bun}_{r,d}\), \(E(n)\) is globally generated. Therefore, we can consider
    \[ H^0(C, E(n)) \twoheadrightarrow E(n). \]
    If we set \(V = \C^N\), where \(N = h^0(E(n))\), then there is a surjection
    \[ V \otimes \mc{O}_C(-n) \twoheadrightarrow E. \]
    Therefore, we can think of \(E\) as a member of a Quot scheme. On the Quot scheme, we can consider the universal family
    \begin{equation*}
    \begin{tikzcd}
        V \otimes \mc{O}_C(-n) \ar[twoheadrightarrow]{r} \ar{d} & \mc{E} \\
        \ms{Quot}_{V \otimes \mc{O}_C(-n), p_E} \times C \ar{r}{p} & \ms{Quot}_{V \otimes \mc{O}_C(-n)}.
    \end{tikzcd}
    \end{equation*}
    Now define the line bundle
    \[ \mc{L} \coloneqq \det R^0 p_* (\mc{E}(m)) \]
    for some \(m \gg 0\). We can also consider the action of \(\mr{SL}(V)\) on the Quot scheme. For some cocharacter \(\lambda\), we can split
    \[ V = V_1 \oplus V_2 \]
    with weights \(s\) and \(N-s\). 
    Then consider
    \[ V_1 \otimes \mc{O}_C(-n) \subset V \otimes \mc{O}_C(-n) \twoheadrightarrow E \]
    and consider the filtration
    \[ 0 \to F^1 E \to E \to E / F^1 E \to 0 \]
    induced by this. Then we see that 
    \[ \lim_{t \to 0} \lambda(t) E = F^1 E \oplus E/F^1 E. \]
    Then we compute the restriction \(\mc{L}_{|\on{gr}_F E}\) and obtain the spaces \(H^0(C, F^1 E(m))\) and \(H^0(C, E(m)/F^1 E(m))\), and therefore the weight of \(\mc{L}\) is given by
    \[ (N-s) p_{F^1 E}(m) - s p_{E/F^1 E}(m), \]
    which recovers GIT.
\end{exm}

We will now consider the unstable locus. Let \(E\) be a vector bundle on \(C\). Then there exists a unique filtration
\[ 0 \subset F^{\ell} E \subseteq \cdots \subseteq F^0 E = E \]
such that
\begin{enumerate}
    \item All \(\on{gr}_F^i E\) are semistable;
    \item We have \(\mu(\on{gr}_F^i E) > \mu(\on{gr}_F^{i-1} E)\) for all \(i\). 
\end{enumerate}
\begin{rmk}
    From the GIT perspective, we should really take \(\mu\) to be the reduced Hilbert polynomial.
\end{rmk}

\subsection{Construction of Theta-stratifications}%
\label{sub:Construction of Theta-stratifications}

Let \(\mf{X}\) be a reasonable algebraic stack.\footnote{Here, reasonable means something like locally of finite presentation with qcqs diagonal.} Now consider some family \(\mc{E}\) over \(\A^1 \times C\). If we restrict to the central fiber, we see that \(\G_m\) acts on \(\on{gr}_F E\), which means we have a map \(\G_m \to \Aut \on{gr}_F E\), and from this we can compute the weight. 

Note that the morphism \(\sigma \colon \ms{Grad}(\mf{X}) \to \ms{Filt}(\mf{X})\) is a \(\Theta\)-contraction, and therefore
\[ \pi_0(\ms{Filt}(\mf{X})) = \pi_0(\ms{Grad}(\mf{X})). \]

\begin{exm}
    Let \(\mf{X} = BG\). Then we have
    \[\ms{Filt}(\mf{X}) = \bigsqcup_{[\lambda]} B P_{\lambda}, \]
    where \(P_{\lambda}\) is the parabolic subgroup corresponding to the cocharacter \(\lambda\), and we can always choose representatives in a single Weyl chamber. Moreover, we have
    \[ \ms{Grad}(\mf{X}) = \bigsqcup_{\lambda \in \mathbb{X}_*(T)/W} B L_{\lambda}, \]
    where \(L_{\lambda}\) is the Levi subgroup corresponding to \(\lambda\).
\end{exm}

\begin{defn}
    A \textit{\(\Theta\)-stratum} in \(\mf{X}\) is a union of connected components \(S \subseteq \ms{Filt}(\mf{X})\) such that
    \[ \ev_1 \colon S \to \mf{X} \]
    is a closed immersion.
\end{defn}

Roughly speaking, closed immersions are injective, so an object in \(S\) consists of an object in \(\mf{X}\) and a ``canonical filtration.''

\begin{defn}
    A \textit{\(\Theta\)-stratification} is an increasing family of open substacks
    \[ \mf{X}_{\leq c} \subseteq \mf{X}, \qquad c \in \Gamma, \]
    where \(\Gamma\) is some totally ordered set, together with \(\Theta\)-strata
    \[ S_c \subseteq \ms{Filt}(\mf{X}_{\leq c}), \]
    which are required to satisfy
    \[ \mf{X}_{\leq c} = \ev_1(S_c) \sqcup \bigcup_{c' < c} \mf{X}_{\leq c'}. \]
    In particular, this implies that
    \[ \mf{X} = \mf{X}^{\mr{ss}} \sqcup \bigsqcup_{c > 0} \ev_1(S_c). \]
\end{defn}

Analogously to the Hilbert-Mumford criterion, we will construct \(\Theta\)-stratifications using numerical invariants.
\begin{defn}
    A \textit{numerical invariant} \(\mu\) with values in a totally ordered set \(\Gamma\) consists of an assignment of a scaling invariant function
    \[ \mu \colon \R^q \setminus \ab\{0\} \to \Gamma \]
    for any \(p \in \mf{X}(\C)\) and any nondegenerate homomorphism \(\gamma \colon \G_m^q \to \Aut(p)\). These are required to satisfy
    \begin{enumerate}
        \item The assignment is invariant under field extensions. This point will not be elaborated on because we are working over \(\C\);
        \item If we have \(\G_m^q \hookrightarrow \G_m^n \hookrightarrow \Aut(p)\), then we have a factorization \(\R^q \setminus \ab\{0\} \to \R^n \setminus \{0\} \to \Gamma\);
        \item If we have some \(\xi \colon S \to \mf{X}\) and \(\gamma \colon \G_{m, S}^q \to \Aut(\xi)\) with finite kernel, the function \(\mu_{\kappa(s)}\) is locally constant on \(S\), where \(s\) ranges over points of \(S\).
    \end{enumerate}
\end{defn}

We will give a naive version of a construction starting with a numerical invariant \(\mu\). Define \(x \in \mf{X}\) to be semistable if \(\mu(f) \leq 0\) for any nondegenerate \(f \in \ms{Filt}(\mf{X})\) satisfying \(f(1) = x\). If \(x\) is unstable, then we define its instability level
\[ M^{\mu}(x) \coloneqq \sup \ab\{ \mu(f) \mid f \in \ms{Filt}(\mf{X}), \ev_1(f) = x\}. \]
If a filtration \(f\) attains this supremum, we call it a ``Harder-Narasimhan filtration.''

\begin{defn}
    A numerical invariant \(\mu\) is \textit{standard} if
    \begin{enumerate}
        \item \(\mu(x)\) and \(\mu(-x)\) are not both positive for \(x \in \R^q \setminus \ab\{0\}\);
        \item Each function \(\mu_{\gamma}\) is strictly quasi-concave, namely for all \(\delta \in [0,1]\) and linearly independent \(x_0, x_1 \in \R^q\) such that \(\mu(x_0), \mu(x_1) \geq 0\), then 
            \[\mu_{\gamma}(\delta x_0 + (1-\delta)x_1) \geq \max\ab\{\mu_{\gamma}(x_0), \mu_{\gamma}(x_1)\}.\] If either value is strictly positive, the inequality is required to be strict.
    \end{enumerate}
\end{defn}

The next condition is that we need some kind of rationality. We consider \(\mathbb{X}_*(T)_{\R}\), and our goal is to find the rational points. This leads us to condition (R): if \(\mu_{\gamma}\) attains a positive value for some \(\gamma\), then it attains a maximal value at a rational point.

\begin{thm}[{\cite{HLinstability}}]
    Let \(\mf{X}\) be a stack and \(\mu\) be a standard numerical invariant satisfying condition (R). Then \(\mu\) defines a \(\Theta\)-stratification if and only if it satisfies the following:
    \begin{itemize}
        \item For any DVR \(R\) with fraction field \(K\) and residue field \(k\), then \(\mu(f_K^{\mr{HN}}) \leq M^{\mu}(f_k)\) (HN-specialization);
        \item For any map \(B \colon T \to \mf{X}\) from a finite type scheme, there exists a quasi-compact open \(\mf{X}' \subseteq \mf{X}\) such that for any point of finite type \(p \in T(k)\) and a filtration \(f\) of \(B(p)\) for which \(\mu(f) > 0\), there exists another filtration \(f'\) of \(B(p)\) with \(\mu(f') \geq \mu(f)\) satisfying \(f'(0) \in \mf{X}'\) (HN-boundedness).
    \end{itemize}
\end{thm}

The point here is that the conditions of being standard and (R) combined with HN-boundedness imply existence and HN-specialization implies uniqueness.



